\section{Step 27: Admit frame count and full-frame period}
Read the four-byte frame count into \texttt{\$s4}. Require a positive value
at most 256. Read the two-byte period into \texttt{\$s5}; require one through
64. Then read and reject nonzero flags. These bounds establish how often the
record walk repeats and which sequence positions require full frames.

Fill \texttt{index\_schedule}. Keeping count and period in saved registers
allows the completed record routine to use temporary registers freely. A count
of 256 remains valid even though the last record index is 255: a count names a
quantity, while a sequence field names one position.

\paragraph{Observe} Prepare \texttt{s27-1} and stop before
\texttt{index\_schedule\_ready}. At \texttt{index\_schedule\_ready}, \texttt{\$s4} is eight,
\texttt{\$s5} is four, and the flags value in \texttt{\$t0} is zero.

\Needspace{7\baselineskip}
\paragraph{Check} What is the last sequence index in a 256-frame file?
\begin{itemize}
\item[A.] 256.
\item[B.] 255.
\item[C.] 257.
\end{itemize}
